\section{Project structure and usage}
\label{app:structure}

\begin{listing}[H]
  \begin{minted}[fontsize=\footnotesize, frame=lines]{text}
notebooks/credit_risk_pipeline.ipynb   Colab front-end (Run all)
modules/creditrisk/
  config.py        paths, dataset URL, column roles, random seed
  data.py          download, rule-based cleaning, stratified split
  eda.py           EDA tables and figures
  transformers.py  grade-wise imputer, IQR clipper, feature engineering
  pipeline.py      PipelineConfig -> scikit-learn / imbalanced-learn Pipeline
  experiments.py   ablation suites, grid search, threshold tuning, test evaluation
  evaluation.py    metrics, cross-validation, ROC / PR / confusion plots
  deep.py          entity-embedding MLP (PyTorch) and TabNet wrapper
  reporting.py     .npy / .h5 feature export, LaTeX tables
  bayesnet/        graph, factor, network, structure, inference, discretize
features/          saved feature sets (.npy) and credit_features.h5
reports/           report.pdf, LaTeX source, figures, generated tables
results/           CSV result of every experiment
tests/             pytest unit tests
  \end{minted}
  \caption{Repository layout.}
  \label{lst:layout}
\end{listing}

\paragraph{Running.}
On Google Colab, open the notebook link from \cref{sec:intro} and choose \emph{Runtime} $\to$ \emph{Run all}.
Locally, install \code{requirements.txt} and PyTorch, then execute the notebook with Jupyter.
The report is compiled with \code{latexmk -pdf -shell-escape report.tex} in \code{reports/latex}.

\section{Unit tests}
\label{app:tests}
The package ships 21 \code{pytest} tests.
They run in a few seconds and cover the following properties.
\begin{itemize}[nosep]
  \item The DAG rejects cycles, and its topological order and Markov blankets are correct.
  \item Variable elimination equals brute-force enumeration on random multi-parent networks, for queries with and without evidence (tolerance $10^{-12}$).
  \item Rejection sampling, likelihood weighting and Gibbs sampling converge to the exact posterior (tolerance 0.01).
  \item Parameter learning recovers the generating CPTs from \num{200000} samples.
  \item A Naive Bayes network gives the same posteriors as scikit-learn's \code{CategoricalNB}.
  \item Missing evidence is marginalized correctly.
  \item TAN and hill climbing produce valid DAGs, and the BIC trace is monotone.
  \item Cleaning, the grade-wise imputer and the IQR clipper behave as specified.
  \item Six representative pipeline configurations fit and predict, and invalid configurations are rejected.
\end{itemize}

\section{Saved feature files}
\label{app:features}
\Cref{tab:features} lists the feature sets written by the notebook.
Every set is stored as separate \code{.npy} arrays and also inside the single file \code{features/credit_features.h5}, with one HDF5 group per set.

\begin{table}[H]
  \centering
  \small
  \begin{tabularx}{\linewidth}{lX}
    \toprule
    \textbf{Folder / HDF5 group} & \textbf{Content} \\
    \midrule
    \code{traditional_logreg} & Output of the tuned logistic-regression preprocessing (imputed, engineered, clipped, one-hot, standardized), with feature names in \code{meta.json} \\
    \code{traditional_<best>} & Output of the preprocessing of the best traditional model \\
    \code{traditional_pca95} & PCA projection retaining 95\% of the variance, with explained-variance ratios \\
    \code{deep_mlp_embeddings} & 64-dimensional penultimate-layer representation of the reference MLP \\
    \bottomrule
  \end{tabularx}
  \caption{Extracted feature sets (\code{X_train}, \code{X_test}, \code{y_train}, \code{y_test} in each).}
  \label{tab:features}
\end{table}
